% 報告会テンプレート（LuaLaTeX対応版）
\documentclass[10pt,a4j,twocolumn]{ltjsarticle}


\usepackage{graphicx}   % \includegraphics を使うためのパッケージ読み込み
\usepackage{imcreport}  % IMC報告会用テンプレートパッケージの読み込み


% ハイフネーションを抑制する。数値が大きいほど、無理にでも抑制しようとする。
\hyphenpenalty=10000\relax
\exhyphenpenalty=10000\relax
\sloppy


\title{中間報告用Lua\LaTeX テンプレートの使い方} % タイトル
\author{発表者氏名}                              % 著者
\studentnumber{学籍番号}                         % 学籍番号
\date{YYYY--MM--DD}                              % 日付
\header{中間報告}                                % 左上のヘッダ内容


\begin{document}

\maketitle

\section{はじめに}

報告会で使われるフォーマットに従った\pLaTeX（またはu\pLaTeX）のテンプレート，``imcreport.sty''を作りました．この資料では作成したフォーマットを使ったサンプルの提示と，テンプレートの使い方の説明をします．

この資料では\pLaTeX のクラスは``jsarticle''を指定していることを前提にしています。資料によってはjarticleを使っていることもありますが、jarticleよりもjsarticleのほうがきれいに組版されるので使える環境ではjsarticleを使うようにしましょう。

\section{図の挿入}

図~\ref{バッテンのラベル}のように、\pLaTeX では図を挿入することができる。
\begin{figure}[bp]
  \centering
  \includegraphics[width=0.45\textwidth]{./sample_figure.pdf}
  \caption{バッテン}
  \label{バッテンのラベル}
\end{figure}

\section{表の挿入}

\pLaTeX で作った表を表~\ref{表のラベル}に示す。
\begin{table}[bp]
  \centering
  \caption{表のタイトル}
  \label{表のラベル}
  \begin{tabular}{c|cccc}
    \hline
          & 項目A & 項目B & 項目C & 項目D \\
    \hline \hline
    項目a & 1     & 2     & 3     & 4 \\
    項目b & 5     & 6     & 7     & 8 \\
    \hline
  \end{tabular}
\end{table}

空間を埋めるためだけの文章です．空間を埋めるためだけの文章です．空間を埋めるためだけの文章です．空間を埋めるためだけの文章です．空間を埋めるためだけの文章です．空間を埋めるためだけの文章です．空間を埋めるためだけの文章です．空間を埋めるためだけの文章です．空間を埋めるためだけの文章です．空間を埋めるためだけの文章です．空間を埋めるためだけの文章です．空間を埋めるためだけの文章です．空間を埋めるためだけの文章です．空間を埋めるためだけの文章です．空間を埋めるためだけの文章です．空間を埋めるためだけの文章です．空間を埋めるためだけの文章です．空間を埋めるためだけの文章です．空間を埋めるためだけの文章です．空間を埋めるためだけの文章です．空間を埋めるためだけの文章です．

\subsection{サブセク}

空間を埋めるためだけの文章です．空間を埋めるためだけの文章です．空間を埋めるためだけの文章です．空間を埋めるためだけの文章です．空間を埋めるためだけの文章です．空間を埋めるためだけの文章です．空間を埋めるためだけの文章です．空間を埋めるためだけの文章です．

\end{document}

